% !PRINTED 113
\weberpage{113}
%<<CONT

\begin{equation}
 x=\sqrt[n]{r}\left(\cos\frac{\varphi}{n}+i\sin\frac{\varphi}{n}\right)
\tag{7}
\end{equation}
eine Grösse, deren $n$te Potenz $=a$ ist, also eine Wurzel der Gleichung
\begin{equation}
 x^n=a,
\tag{8}
\end{equation}
wenn unter $n$ eine beliebige positive ganze Zahl verstanden wird. Derselben Forderung genügt aber auch jeder der Werthe
\begin{equation}
 x_k=\sqrt[n]{r}\left[\cos\left(\frac{\varphi+2k\pi}{n}\right)+i\sin\left(\frac{\varphi+2k\pi}{n}\right)\right],
\tag{9}
\end{equation}
wenn $k$ eine beliebige ganze Zahl ist. In (9) sind aber $n$ und nicht mehr verschiedene Werthe enthalten, die man erhält, wenn man
\begin{equation}
 k=0,1,2,\ldots,n-1
\tag{10}
\end{equation}
setzt; denn vermehrt man $k$ um ein Vielfaches von $n$, so ändert (9) seinen Werth nicht, während die $n$ Werthe (10) lauter verschiedene Werthe von $x$ ergeben, weil sie verschiedene Phasen haben.

Die Gleichung (8) hat demnach nicht nur eine, sondern $n$ verschiedene Wurzeln. Die geometrischen Bilder der Werthe $x_k$ liegen alle auf einem Kreise mit dem Radius $\sqrt[n]{r}$, und zwar um den Winkel $\frac{2\pi}{n}$ von einander entfernt. Man erhält den ersten dieser Punkte dadurch, dass man den gegebenen Winkel $\varphi$ in $n$ gleiche Theile theilt. Der Radius $\sqrt[n]{r}$ ist grösser oder kleiner als $r$, je nachdem $r$ kleiner oder grösser als $1$ ist. Die beistehende Fig. 1 zeigt diese Verhältnisse für $n=3$ unter der Voraussetzung, dass $r>1$ ist.


% Exact 300-dpi crop from the registered Volume-I scan, printed Fig. 1.
\weberfigureimage{weber-b1-fig01.png}{0.50}


Man kann die verschiedenen Werthe von $x_k$ dadurch erhalten, dass man den ersten von ihnen
\[
 x_0=\sqrt[n]{r}\left(\cos\frac{\varphi}{n}+i\sin\frac{\varphi}{n}\right)
\]
mit den verschiedenen Werthen von
%CONT>>
